\chapter{Globale Assembly, Elektroden und Auswertung}\label{chap:algebra}
\section{Freiheitsgrade und Materialregionen}
Ein skalares Feld trägt einen DOF pro aktivem Knoten, Mechanik drei in
knotenweiser xyz-Reihenfolge. Piezo verwendet global $[\mathbf u,\boldsymbol\phi]$;
die allgemeine Feldreihenfolge bei drei Feldern ist $[\mathbf u,T,\phi]$.
Aktive Knoten sind die sortierten in Volumenzellen benutzten Knoten.
Netzknotenindizes, kompakte DOF-Indizes und externe Entitäts-IDs sind
verschiedene Mengen. Ein ungenutzter Netzpunkt erzeugt keinen Solver-DOF;
Ergebnisarrays müssen über \texttt{node\_ids} zurückabgebildet werden.

Mit der Extraktionsmatrix $\mathbf R_e$ gilt
\begin{equation}\label{eq:assembly}
\mathbf A=\sum_e\mathbf R_e^T\mathbf A_e\mathbf R_e,
\qquad\mathbf b=\sum_e\mathbf R_e^T\mathbf b_e+\mathbf b_{\rm boundary}.
\end{equation}
COO-Einträge werden addiert und nach CSR/CSC überführt; gemeinsame Knoten
summieren Beiträge. Jede aktive Zelle benötigt je Physik genau eine
Materialzuordnung. Identische überlappende Zuordnungen sind zulässig,
widersprüchliche Zuordnungen, Lücken und ungenutzte Definitionen nicht.
Perfekte konforme Interfaces teilen Primär-DOFs; damit sind $T,\phi,\mathbf u$
stetig, während Gradienten und Rohspannungen springen können.
Die schwache Form erzwingt Interface-Fluss-/Traktionsgleichgewicht im
Galerkin-Sinn, nicht punktweise Kontinuität jeder exportierten Probe.
\index{Assembly}\index{Materialregion}

\section{Symmetrieerhaltende Dirichlet-Elimination}
$\mathbf P_f$ sei der diagonale Projektor auf freie DOFs, $\mathbf g$ habe
vorgegebene Werte auf festen und null auf freien DOFs. Für eine reelle
symmetrische Matrix ist die implementierte Vollraumform
\begin{equation}\label{eq:dirichlet}
\widetilde{\mathbf A}=\mathbf P_f\mathbf A\mathbf P_f+(\mathbf I-\mathbf P_f),
\qquad\widetilde{\mathbf b}=\mathbf P_f(\mathbf b-\mathbf A\mathbf g)+\mathbf g.
\end{equation}
Die freien Gleichungen sind äquivalent zu
$\mathbf A_{ff}\mathbf x_f=\mathbf b_f-\mathbf A_{fc}\mathbf g_c$.
Symmetrie bleibt erhalten; ursprüngliche $\mathbf A,\mathbf b$ werden für
Reaktionen bewahrt. Die Einheitszeilen sind algebraische Vorgaben und besitzen
keine ursprüngliche physikalische Steifigkeitseinheit.
Identische Mehrfachvorgaben sind erlaubt, widersprüchliche Werte Fehler.

Eine Zwei-Feder-Kette liefert als Handrechnung
\begin{equation}\label{eq:spring-example}
\mathbf A=\begin{pmatrix}1&-1&0\\-1&2&-1\\0&-1&1\end{pmatrix},
\quad\mathbf b=(0,0,3)^T,\quad x_0=2.
\end{equation}
Die freien Gleichungen geben $x_1=5$, $x_2=8$; das ursprüngliche Residuum
ist $(-3,0,0)^T$. Die feste Einheitszeile hätte Residuum null und würde die
Lagerreaktion verlieren. Zahlen gelten für Federsteifigkeit 1 N/m,
Verschiebungen in m, Endlast 3 N; der Test ist algebraisch und kein
Kleine-Verzerrungs-Experiment.
\index{Dirichlet-Elimination}\index{Reaktion}

\section{Equipotentiale Elektroden und Gesamtladung}
Eine Inzidenzmatrix $\mathbf P_e$ bildet alle Potentialknoten einer schwebenden
Elektrode auf genau einen Leiter-DOF ab. Mechanik und andere elektrische
DOFs bleiben unverändert. Mit $\mathbf x=\mathbf P_e\mathbf y$ gilt
\begin{equation}\label{eq:electrode-condensation}
\mathbf A_r=\mathbf P_e^T\mathbf A\mathbf P_e,\qquad
\mathbf b_r=\mathbf P_e^T\mathbf b-\mathbf q_L.
\end{equation}
$\mathbf q_L$ trägt die vorgegebene Gesamtleiterladung auf dem Leiter-DOF.
Nach Rückabbildung $\mathbf r=\mathbf A\mathbf x-\mathbf b$ folgt
$\sum_{i\in L}r_i=-Q_L$. Einzelne Elektrodenknoten dürfen Residuen besitzen;
ihre Summe ist die freie elektrische Gleichung. Daher werden sie in der
Auswertung nicht fälschlich wie voneinander unabhängige freie Potentiale
behandelt. Alle sechs Triangle6-Knoten, auch Eckknoten mit verschwindender
konstanter Last, gehören zur Elektrode.

Die Kondensation erhält Symmetrie. Nachfolgende Dirichlet-Vorgaben fixieren
feste Leiterpotentiale und den erforderlichen Bezug. Sich an einem Knoten
berührende unterschiedliche Elektroden, doppelte Facetten, zusätzliche
Potentialvorgaben auf Floating-Knoten und simultane D-Randflussvorgaben auf
Leiterflächen sind widersprüchliche Daten. Der Solver lehnt sie ab.
\index{Schwebende Elektrode}

\subsection{Handrechnung mit vorgegebener Leiterladung}
Ein 1-m-Quader, Fläche 1 m$^2$, $e=0{,}5$ C/m$^2$, $\kappa=2$ F/m und
vorgegebene Dehnung $\varepsilon=0{,}1$ besitzt
$D=0{,}05-2\phi'$, $Q_{L,\rm rechts}=-D$.
Leerlauf $Q_L=0$ gibt $\phi(1)=0{,}025$ V; Kurzschluss gibt $D=0{,}05$
C/m$^2$ und $Q_L=-0{,}05$ C. Mit $Q_L=1{,}95$ C ist
$D=-1{,}95$ C/m$^2$ und $\phi(1)=1$ V. Das Vorzeichen der Leiterladung
ist damit direkt nachrechenbar, ohne einen Solver als Referenz zu verwenden.

\section{Feldrekonstruktion, Energien und Bilanzen}
Primärfelder werden mit den ursprünglichen Formfunktionen interpoliert;
Messpunkte sind keine Nächster-Knoten-Auswahl. Tet4 liefert konstante rohe
Gradienten, Spannung, Fluss, E und D. Tet10 liefert rohe Schwerpunktproben
mit explizitem Auswerteort. Auf konformen Netzen bleiben die Primärfelder
stetig, die abgeleiteten Werte im Allgemeinen unstetig. Nodale Mittelung,
Glättung oder Projektion darf nicht als Verbesserung der physikalischen
Genauigkeit interpretiert werden.

Zellenergie/-enthalpie ist bei Tet10 ein Quadraturintegral; Schwerpunktprobe
mal Volumen ist im Allgemeinen falsch. Für $T=x^2$ in
Abschnitt~\ref{sec:unit-tet} wäre die Schwerpunkt-Energiedichte $8(1/4)^2$
und ihr Volumenprodukt $1/12$ W K statt $2/15$ W K.
Tet4 verwendet die dokumentierte thermische Zellmittelkonvention.
Komplexe Felder werden vor der Betrags-/Phasenbildung interpoliert;
$|\sum N_a\hat u_a|$ ist nicht $\sum N_a|\hat u_a|$.
Nullgefüllte inaktive Daten werden mit Gültigkeitsmasken von echten Nullwerten
unterschieden.

Freie Residuen, Reaktionen, Leiterladungen, Wärmeleistung, Kraft und Moment
werden am ursprünglichen System rekonstruiert. Die Globalkontrollen aus
Gleichungen~\ref{eq:thermal-balance} und \ref{eq:charge-balance} bleiben
getrennt; ihre SI-Einheiten dürfen nicht in eine einzige unskalierte Norm
vermengt werden. Eine Verletzung unterhalb eines numerisch begründeten
Rundungsbudgets ist anders zu beurteilen als ein fehlender Bilanznachweis.

\subsection*{Implementierung und Nachweise.}
{\small\raggedright
\path{src/fem_toolbox/assembly/__init__.py},
\path{src/fem_toolbox/boundary/electrodes.py},
\path{src/fem_toolbox/solvers/piezoelectric.py};
\path{tests/unit/test_assembly.py},
\path{tests/patch/test_electrodes.py},
\path{tests/patch/test_material_interfaces.py},
\path{tests/unit/test_tet10_piezoelectric_integration.py}.
Die Quellen-/Nachweismatrix benennt die weitere Topologie- und Exportabdeckung.

\par}
